\section{Verifiable Pseudorandom Secret Sharing}


We now introduce an extension to pseudorandom secret sharing that we call a  \emph{Verifiable Pseodurandom Secret Sharing} (VPSS) scheme.
We begin by motivating the need for a VPSS,
and then formally define VPSS and its security properties.
Finally,
in Section~\ref{dvrf-concrete},
we give a concrete VPSS construction, \DVRFOne, that we later use as a building block for \glitter.


\subsection{Motivation}

A verifiable random function (VRF)~\cite{MicaliRV99} is a keyed PRF,
such that the PRF can be evaluated using only knowledge of a secret key,
but is publicly verifiable given a public key.
In particular, in addition to outputting the evaluation of the VRF,
it also outputs a zero-knowledge proof that the VRF was evaluated correctly.
A \emph{distributed} VRF~\cite{NaorPR99,Dodis03,GalindoLOW21} allows the evaluation algorithm to be partitioned among a set of participants,
all whom are equally trusted.

%The efficiency of a VRF depends on its domain and co-domain.
%For example,
%for pairing-based VRFs~\cite{Dodis03},
%the evaluation of the pseodorandom function outputs an element in the target group,
%and so verifying that the operation was done correctly can be done using only a small number of group operations.
%
However, the use case of generating deterministic nonces for threshold Schnorr signatures presents a slightly different challenge.
Instead of directly verifying that the nonce $r_i$ was correctly generated,
we instead need to verify correctness in zero-knowledge,
with respect to a commitment $R_i = g^{r_i}$.
One approach in the literature is to employ non-algebraic pseudorandom functions to generate the nonce,
and then prove in zero-knowledge the correctness of the corresponding commitment~\cite{NickRSW20,GarillotKMN21}.
Unfortunately, generating such zero-knowledge proofs requires higher computational and complexity overhead.

We instead take a different approach towards verifying that a distributed pseudorandom function was correctly performed by a set of parties.
Instead of each participant outputting a zero-knowledge proof that they individually performed the action correctly,
we observe that the correctness of the evaluation can instead be \emph{collectively} publicly verified,
assuming some threshold of honest participants.
%
Such an observation leads naturally to employing a pseudorandom secret sharing scheme~\cite{CramerDI05},
which is categorized by all parties holding a set of secret key shares,
and individually and non-interactively being able to generate secret shares
of additional pseudorandom values.
We show that pseudorandom secret sharing schemes has a natural extension to the publicly verifiability setting.
Finally,
we give a concrete and efficient VPSS scheme,
where all participants can publish commitments to their generated shares,
and perform polynomial interpolation over elements in the public domain to ensure correctness.

We next build upon these observations by formalizing the notion of a verifiable pseudorandom secret sharing scheme.

\subsection{VPSS Definition and Notions of Security}

We now present an extension on pseudorandom secret sharing that we call a
\emph{Verifiable Pseudorandom Secret Sharing} (VPSS) scheme.
In particular,
a VPSS defines an additional verify algorithm to ensure that the pseudorandom function was performed correctly by collectively verifying outputs by all participants.
%in a coalition $\coalition$,
%where $\coalition \subseteq \set{n}$.
%In doing so,
%we can define a more lightweight scheme than is possible in the setting where each partial evaluation is individually verifiable.

%We now more formally define a VPSS,
%and its required security properties.

\begin{definition}
A \defn{Verifiable Pseudorandom Secret Sharing} scheme is the tuple of algorithms
$\dvrf = (\setup, \keygen, \generate, \verify, \combinepub, \recover)$,
such that:

\begin{itemize}
  \item $\setup(\usecp)$: Accepts as input the security parameter $\secp$,
and outputs public parameters $\pp$,
where $\pp$ is given as implicit input to all other algorithms.
  \item $\keygen(n, t,\minparticipants) \rightarrow \bot/(\SK_1, \ldots, \SK_n)$:
    A probabilistic algorithm that accepts as input the total number of participants $n$,
    a corruption threshold $t$,
    and the minimum number of parties $\minparticipants$ required to participate in $\generate$.
    On failure,
    output $\bot$,
    otherwise,
    output $n$ secret keys, one for each of the $n$ participants.
  \item $\generate(\partyid, \SK_\partyid, \vrfin) \rightarrow
    (\vrfprivout_\partyid, \vrfpubout_\partyid)$:
    A deterministic algorithm that accepts as input a participant identifier $\partyid$,
the secret key for that participant $\SK_\partyid$,
and some input $\vrfin$.
Generates the pseudorandom secret share $\vrfprivout_\partyid \in \vrfprivdomain$ from
some domain $\vrfprivdomain$,
using $\SK_\partyid$ as the random seed and $\vrfin$ as the corresponding input.
Then, generates its public commitment $\vrfpubout_k \in \vrfpubdomain$
to $\vrfprivout_\partyid$ from some domain $\vrfpubdomain$.
Outputs $(\vrfprivout_\partyid, \vrfpubout_\partyid)$.
%  \item $\commit(\vrfprivout) \rightarrow \vrfprivout$
%Accepts as input a VRF evaluation $\vrfprivout \in \vrfprivdomain$
%and outputs a public commitment $\vrfpubout \in \vrfpubdomain$ from some domain $\vrfpubdomain$.
  \item $\verify(t, \minparticipants, \coalition, \{ \vrfpubout_j \}_{j \in \coalition}) \rightarrow \zo$:
A deterministic algorithm that accepts as input the corruption threshold $t$,
the minimum number of participants $\minparticipants$,
a coalition of participants $\coalition$ such that $\lvert \coalition \rvert \geq \minparticipants$,
and a set of commitments to pseudorandom secret shares $( \vrfpubout_j )_{j \in \coalition}$  for that coalition.
Outputs $1$ to indicate that the secret sharing was correctly performed,
otherwise, output $0$.
%\ian{Is it important that $\verify$ and $\combinepub$ be separate functions? Wouldn't it be simpler for $\combinepub$ to just output $\bot$ if the inputs were inconsistent, and $\vrfpubout$ otherwise? You use $\verify$ a lot in the paper, though, so changing it at this point is probably not a good idea.}
%\chelsea{I think it is conceptually clearer with these as seperate functions}
%  \item $\combinepriv(t, \minparticipants, \coalition, \{ \vrfprivout_j \}_{j \in \coalition}) \rightarrow \vrfprivout $:
%Accepts as input the corruption threshold $t$, the minimum number of participants $\minparticipants$,
%a coalition of participants $\coalition$ such that $\lvert \coalition \rvert \geq \minparticipants$,
%and the set of VRF partial evaluations.
%Outputs the aggregated VRF evaluation $\vrfprivout$.
  \item $\combinepub(t, \minparticipants, \coalition, \{ \vrfpubout_j \}_{j \in \coalition}) \rightarrow \vrfpubout$:
A deterministic algorithm that accepts as input the corruption threshold $t$, the minimum number of participants $\minparticipants$,
a coalition of participants $\coalition$ such that $\lvert \coalition \rvert \geq \minparticipants$,
and the set of commitments to pseudorandom secret shares.
Outputs the commitment to the aggregated pseudorandom secret $\vrfpubout$.
    \item $\recover(t, \minparticipants, \coalition, \{ \vrfprivout_j \}_{ j \in \coalition}) \rightarrow \bot /  (\vrfprivout, \vrfpubout )$:
      A deterministic algorithm that accepts a corruption threshold $\thresh$,
      the minimum number of participants $\minparticipants$,
      a coalition $\coalition \subseteq \set{n}, \lvert \coalition \rvert \geq \minparticipants$,
      and a set of shares  $\{ \vrfprivout_j \}_{j \in \coalition}$.
      It outputs $\bot$ if $\coalition \not\subseteq \set{n}$,
      $\lvert \coalition \rvert < \minparticipants$,
      or if the shares are inconsistent.
      Otherwise, it recovers $\vrfprivout$ using the set of shares,
      derives the corresponding commitment $\vrfpubout$, and
      outputs $(\vrfprivout, \vrfpubout)$.

\end{itemize}
\end{definition}


\paragraph{Correctness.}
A VPSS is \defn{correct} if for all $\secp$,
possible inputs $\vrfin$ and choices of  $t, \minparticipants, n \in \mathbb{N}$
where $t \le \minparticipants \le n$,
%\chelsea{allowable is defined by the construction, we cannot define what is
%allowable here. E.e, we shouldn't enforce honest majority at this level.}
%and all $\coalition  \subseteq \set{n}$ with $|\coalition| \ge \minparticipants, $,
when $\pp \gets \dvrf.\setup(\usecp)$
and $(\SK_1, \ldots, \SK_n) \rgets \dvrf.\keygen(n, t, \minparticipants)$,
there exists $\vrfprivout \in \vrfprivdomain, \vrfpubout \in \vrfpubdomain$ such that Equation~\ref{eq:vrf-correct} holds:

\begin{align} \label{eq:vrf-correct}
  \begin{split}
    \text{ for all } i \in \set{n} \text{ and for all } \coalition_\ell \subseteq \set{n}, \lvert \coalition_\ell \rvert \geq \minsigners,
    \text{ when } \\
    \dvrf.\generate(i, \SK_i, \vrfin) \rightarrow (\vrfprivout_i, \vrfpubout_i),
    \text{ then } \\
    \dvrf.\verify(t, \minparticipants, \coalition_\ell, \{ \vrfpubout_j \}_{j \in \coalition}) = 1,
    \text{ and } \\
    \dvrf.\recover(t, \minparticipants, \coalition_\ell, (\vrfprivout_j)_{j \in \coalition})
    \rightarrow (\vrfprivout, \vrfpubout), \text{ and} \\
    \dvrf.\combinepub(t, \minparticipants, \coalition_\ell, (\vrfpubout_j)_{j \in \coalition})  \rightarrow \vrfpubout
  \end{split}
\end{align}

% TODO cite NaorPR99,Dodis03
\subsubsection{Security.}
Similarly to verifiable random functions~\cite{MicaliRV99,NaorPR99,Dodis03,GalindoLOW21},
we say that a VPSS is \emph{secure} if it is unique, verifiable, and pseudorandom.
%\chelsea{These notions do not come from the PSS paper, they come from the dodis paper. The PSS paper never defined these.}

\begin{figure}[t]
  \centering
	%\begin{linedgamespec}
	\begin{pchstack}[boxed]
		\begin{pcvstack}
		\procedure{ $\UniqExp_{\adv,\dvrf}(\secp, n, t, \minparticipants)$ }{
      \pp \gets \dvrf.\setup(\usecp) \\
      (\corrupt, \advstate) \rgets \adv(\pp, n, t, \minparticipants) \\
      \pcreturn \bot\ \pcif \lvert \corrupt \rvert \geq t \\
      \honest \gets \set{n} \setminus \corrupt \\
      (\SK_1, \ldots, \SK_n) \rgets \dvrf.\keygen(n, t, \minparticipants) \\
      (\vrfin, \mathsf{out}) \rgets \adv^{\oracle^{\generate}}((\SK_j)_{j \in \corrupt}, \advstate) \\
      (\coalition, (\vrfpubout_i)_{i \in \coalition \cap \corrupt}),
      (\coalition', (\vrfpubout_i')_{i \in \coalition' \cap \corrupt} )  \gets \mathsf{out} \\
      (\vrfprivout_j, \vrfpubout_j) \gets \dvrf.\generate(j, \SK_j, \vrfin), j \in \honest \\
      S_1 \gets ( \vrfpubout_i )_{i \in \coalition \cap \corrupt} \cup ( \vrfpubout_i )_{i \in \coalition \cap \honest}  \\
      S_2 \gets ( \vrfpubout_i' )_{i \in \coalition' \cap \corrupt} \cup ( \vrfpubout_i )_{i \in \coalition' \cap \honest}  \\
      \pcreturn 0\ \pcif \dvrf.\verify(t, \minparticipants, \coalition, S_1) \neq 1 \\
      \pcreturn 0\ \pcif \dvrf.\verify(t, \minparticipants, \coalition', S_2) \neq 1 \\
      %S_1' \gets ( \vrfprivout_i )_{i \in \coalition \cap \corrupt} \cup ( \vrfprivout_i )_{i \in \coalition \cap \honest}  \\
      %S_2' \gets ( \vrfprivout_i' )_{i \in \coalition' \cap \corrupt} \cup ( \vrfprivout_i )_{i \in \coalition' \cap \honest}  \\
      %\pcif \dvrf.\combinepub(t, \minparticipants, \coalition, S_1) = \bot \\
      %\textbf{or }  \dvrf.\combinepub(t, \minparticipants,  \coalition', S_2) = \bot \\
      %\pcind \pcreturn 0 \\
      %\text{\ian{Now that you're using $\combinepub$,}}\\
      %\text{\ian{the above three lines can be removed.}}\\
      \vrfpubout \gets \dvrf.\combinepub(t, \minparticipants, \coalition, S_1) \\
      \vrfpubout' \gets \dvrf.\combinepub(t, \minparticipants, \coalition', S_2) \\
      \pcif \vrfpubout \neq \vrfpubout' \\
      %\pcif \dvrf.\combinepub(t, \minparticipants, \coalition, S_1) \neq \dvrf.\combinepub(t, \minparticipants,  \coalition', S_2) \\
      \pcind \pcreturn 1 \\
      \pcreturn 0
		}
		\end{pcvstack}
    		\pchspace
		\begin{pcvstack}
			\procedure{ $\BigO^{\generate}(\partyid, \vrfin_i)$ }{
        \pccomment{$\partyid$ denotes the participant identifier} \\
        (\vrfprivout_\partyid, \vrfpubout_\partyid) \gets \dvrf.\generate(\partyid, \SK_\partyid, \vrfin_i) \\
        \pcreturn (\vrfprivout_\partyid, \vrfpubout_\partyid)
			}
		\end{pcvstack}
	\end{pchstack}
	%\end{linedgamespec}
\caption{Uniqueness game for a VPSS.
  }
\label{fig:dcvrf-unique}
\end{figure}

%
\paragraph{Uniqueness.}
Intuitively,
a VPSS is unique if it produces exactly one commitment to a (pseudorandom) value for each input $\vrfin$,
regardless of the choice of coalition.
%\ian{You had the next sentence backwards.  For the same input $\vrfin$, different coalitions $\coalition$ \emph{must} give the same commitment $\vrfpubout$. (In the game, the adversary \emph{wins} if the $\vrfpubout$ values are different. You had it correct a couple of paragraphs down from here.) Currently in Arctic, different coalitions will result in a different $\vrfin$, but that's not what we're talking about here.}
More specifically,
picking different coalitions should always result in the same commitment $\vrfpubout$,
when the input $\vrfin$ remains the same.
We show the VPSS uniqueness experiment in Figure~\ref{fig:dcvrf-unique}.

In the uniqueness experiment,
the adversary is allowed to query honest participants for shares and commitments on inputs of its choosing.
The adversary then outputs
the evaluation input $\vrfin$,
two coalitions $\coalition, \coalition'$,
and two sets of commitments
$(\vrfpubout_i)_{i \in \coalition \cap \corrupt}$,
$( \vrfpubout_i')_{i \in \coalition' \cap \corrupt}$.

The environment then derives the honest players' shares and commitments on $\vrfin$,
producing
$(\vrfprivout_j, \vrfpubout_j)_{j \in \honest}$.
After deriving the sets
$S_1 \gets ( \vrfpubout_i )_{i \in \coalition \cap \corrupt} \cup ( \vrfpubout_i )_{i \in \coalition \cap \honest}$
and $S_2 \gets ( \vrfpubout_i' )_{i \in \coalition' \cap \corrupt} \cup ( \vrfpubout_i )_{i \in \coalition' \cap \honest}$,
the adversary loses if the verification algorithm outputs 0 on either set.
%or if $\dvrf.\combine$ outputs $\bot$ for either set.
%
If both sets verify,
%and the combine algorithm succeeds for both sets,
the adversary wins if the corresponding commitments for each set are not equal.

We define uniqueness for a VPSS more formally in Definition~\ref{def:dcvrf-unique}.

\begin{definition} \label{def:dcvrf-unique}
Let the advantage of an adversary $\adv$ playing the uniqueness game as defined in Figure~\ref{fig:dcvrf-unique} be as follows:

\[
  \advant{unq}{\adv,\dvrf}(\secp, n, t, \minparticipants)
  = \bigabs{\Pr[ \UniqExp_{\adv,\dvrf}(\secp, n, t, \minparticipants)= 1]}
\]


  A VPSS $\dvrf$ is \defn{unique} if for all PPT adversaries $\adv$,
  $\advant{unq}{\adv,\dvrf}$ is a negligible function of $\secp$,
  for $n, t, \minparticipants \in \mathbb{N}$,
  $t \leq \minparticipants \leq n$.
\end{definition}




\begin{figure}[t]
  \centering
	\begin{pchstack}[boxed]
		\begin{pcvstack}
		\procedure{ $\VerifExp_{\adv,\dvrf}(\secp, n, t, \minparticipants)$ }{
      \pp \gets \dvrf.\setup(\usecp) \\
      (\corrupt, \advstate) \rgets \adv(\pp, n, t, \minparticipants) \\
      \pcreturn \bot\ \pcif \lvert \corrupt \rvert \geq t \\
      \honest \gets \set{n} \setminus \corrupt \\
      (\SK_1, \ldots, \SK_n) \rgets \dvrf.\keygen(n, t, \minparticipants) \\
      (\vrfin, \coalition, (\vrfpubout_j)_{j \in \coalition \cap \corrupt}) \rgets \adv^{\oracle^{\generate}}((\SK_j)_{j \in \corrupt}, \advstate) \\
      \pcfor j \in \coalition \pcdo \\
      \pcind (\vrfprivout_j', \vrfpubout_j') \gets \dvrf.\generate(j, \SK_j, \vrfin)  \\
      \pcif (\vrfpubout_j)_{j \in \coalition \cap \corrupt} = (\vrfpubout_j')_{j \in \coalition \cap \corrupt} \\
      \pcind \pcreturn 0 \\
      \text{\algcom $\adv$ must deviate from the protocol to win} \\
      S \gets (\vrfpubout_k)_{k \in \coalition \cap \corrupt} \cup (\vrfpubout_j')_{j \in \coalition \cap \honest} \\
      \pcif \dvrf.\verify(t, \minparticipants, \coalition, S) = 1 \\
      \pcind \pcreturn 1 \\
      \pcreturn 0
		}
		\end{pcvstack}
    		\pchspace
		\begin{pcvstack}
			\procedure{ $\BigO^{\generate}(\partyid, \vrfin_i)$ }{
        \pccomment{$\partyid$ denotes the participant id} \\
        (\vrfprivout_\partyid, \vrfpubout_\partyid) \gets \dvrf.\generate(\partyid, \SK_\partyid, \vrfin_i) \\
        \pcreturn (\vrfprivout_\partyid, \vrfpubout_\partyid)
			}
		\end{pcvstack}
	\end{pchstack}
	%\end{linedgamespec}
\caption{Verifiability game for a VPSS.
%\ian{You need to check that $C$ has size at least $\minparticipants$ (return 0 if not).}
%\chelsea{This needs to be done by verify, not the environment}
%\ian{$\honest$ should be $\coalition \setminus \corrupt$, and defined after $\adv$ returns $\coalition$.}
%\chelsea{honest is defined for all possible parties, not just the coalition}
  }
\label{fig:dcvrf-verif}
\end{figure}

\paragraph{Verifiability.}
Intuitively,
a VPSS is verifiable if given a set of commitments to shares from a coalition of participants,
the verify algorithm will detect if some subset of players deviated from the protocol.
We show the VPSS verifiability experiment in Figure~\ref{fig:dcvrf-verif}.

In the experiment,
the adversary is allowed to query honest participants for shares and commitments on inputs of its choosing.
The adversary then outputs a coalition $\coalition$,
a VPSS input $\vrfin$,
and a set of commitments
$(\vrfpubout_i)_{i \in \coalition \cap \corrupt}$.
The environment then follows the protocol, deriving both the corrupt and honest players' commitments on $\vrfin$,
and producing
$(\vrfpubout_j')_{j \in \coalition}$.
The adversary loses if its output is identical to the honestly derived commitments for the corrupted players.
%
Then,
the environment checks if the set of commitments
$(\vrfpubout_i)_{i \in \coalition \cap \corrupt} \cup (\vrfpubout_j')_{j \in \coalition \cap \honest}$ are valid with respect to $\coalition$ and $\vrfin$.
If so,
the adversary wins,
otherwise,
it loses.

We define verifiability for a VPSS more formally in Definition~\ref{def:verifiable}.

\begin{definition} \label{def:verifiable}
Let the advantage of an adversary $\adv$ playing the verifiability game as
defined in Figure~\ref{fig:dcvrf-verif} be as follows:

\[
  \advant{verf}{\adv,\dvrf}(\secp, n, t, \minparticipants)
  = \bigabs{\Pr[ \VerifExp_{\adv,\dvrf}(\secp, n, t, \minparticipants)= 1]}
\]

  A VPSS $\dvrf$ is \defn{verifiable} if for all PPT adversaries $\adv$,
  $\advant{verf}{\adv,\dvrf}$ is a negligible function of $\secp$,
  for $n, t, \minparticipants \in \mathbb{N}$,
  $t \leq \minparticipants \leq n$.
\end{definition}

\begin{figure}[t]
  \centering
	\begin{pchstack}[boxed]
		\begin{pcvstack}
		\procedure{ $\PseudoExp_{\adv,\dvrf}(\secp, n, t, \minparticipants)$ }{
     b \rgets \zo \\
     z \rgets \zo \\
      \text{\algcom random coin for trivial losing conditions}\\
%     z \rgets \zo \text{\algcom random coin for losing conditions} \\
     Q \gets \emptyset \\
     \pp \gets \dvrf.\setup(\usecp) \\
      (\corrupt, \advstate) \rgets \adv(\pp, n, t, \minparticipants) \\
      \pcreturn \bot\ \pcif \lvert \corrupt \rvert \geq t \\
     \honest \gets \set{n} \setminus \corrupt \\
     (\SK_1, \ldots, \SK_n) \rgets \dvrf.\keygen(n, t, \minparticipants) \\
      \mathsf{out} \rgets \adv^{\oracle^{\generate}}((\SK_j)_{j \in \corrupt}, \advstate) \\
      (b', \vrfin^*, \coalition^*, \{ \vrfpubout_j \}_{j \in \corrupt}) \gets \mathsf{out} \\
      \pcif b=1  \\
      \pcind \pcreturn z\ \pcif Q[\vrfin^*]  = \bot \\
      \pcind (\vrfpubout, \{ \vrfpubout_i\}_{i \in \honest}) \gets Q[\vrfin^*] \\
      \pcind \pcif \dvrf.\verify(t, \minparticipants, \coalition^*, \{ \vrfpubout_k \}_{k \in \coalition^*}) \neq 1 \\
      \pcind \pcind  \pcreturn z \\
      \pcind \pcif \dvrf.\combinepub(t, \minparticipants, \coalition^*, \{ \vrfpubout_k \}_{k \in \coalition^*}) \neq \vrfpubout \\
      \pcind \pcind  \pcreturn 1 \\
      \pcreturn 1\ \pcif b \iseq b' \\
      \pcreturn 0
		}
		\end{pcvstack}
    		\pchspace
		\begin{pcvstack}
			\procedure{ $\BigO^{\generate}(\partyid, \vrfin)$ }{
        \pccomment{$\partyid$ denotes the participant identifier} \\
      \pcreturn \bot\ \pcif \partyid \not\in \honest \\
     %   Q \gets Q \cup \{ \vrfin_i \} \\
     \pcif b=0 \\
      \pcind (\vrfprivout_\partyid, \vrfpubout_\partyid) \gets \dvrf.\generate(\partyid, \SK_\partyid, \vrfin) \\
      \pcind \pcreturn (\vrfprivout_\partyid, \vrfpubout_\partyid)  \\
      %
      \text{\algcom $b=1$ case} \\
      \pcif  Q[\vrfin] \neq \bot \\
      \pcind  (\vrfpubout, \{ (\vrfprivout_i, \vrfpubout_i) \}_{i \in \honest} ) \gets Q[\vrfin] \\
      %
      \pcelse \\
       \pcind \vrfprivout' \rgets  \vrfprivdomain \\
       \pcind  \mathsf{in} \gets(\vrfin, \vrfprivout', \corrupt, n, (\sk_j)_{j \in \corrupt})   \\
       \pcind (\vrfpubout, \{ (\vrfprivout_i, \vrfpubout_i ) \}_{i \in \honest} ) \gets \SIMEval(\mathsf{in}) \\
       \pcind Q[\vrfin] = (\vrfpubout, \{ (\vrfprivout_i, \vrfpubout_i ) \}_{i \in \honest}) \\
    \text{ \algcom $\SIMEval$ is a simulating algorithm} \\
    \text{ \algcom defined by the VPSS construction.} \\
        \pcreturn (\vrfprivout_\partyid, \vrfpubout_\partyid)
        % the adversary shouldn't be allowed to query this oracle for its
        % challenge, but it should be allowed to learn the public honest commitments?
			}
		\end{pcvstack}
	\end{pchstack}
	%\end{linedgamespec}
\caption{Pseudorandomness game for a VPSS.
  }
\label{fig:dcvrf-pseudo}
\end{figure}

\paragraph{Pseudorandom.}
Intuitively,
a VPSS is pseudorandom if the adversary has negligible advantage distinguishing between a real VPSS output from one that is randomly sampled.
We show the VPSS pseudorandomness experiment in Figure~\ref{fig:dcvrf-pseudo}.

In the pseudorandomness experiment,
the environment begins by picking a random bit $b \rgets \zo$.
It then performs key generation,
and initializes the adversary with the public parameters and the corrupted parties secret key shares.

The adversary is allowed to query $\oracle^{\generate}$ on honest participants for shares and commitments on inputs of the adversary's choosing.
If $b=0$,
the environment responds by following the VPSS protocol.
If $b=1$,
the environment first checks to see if it has responded to the query before,
replying with the response if so.
Otherwise, it uses a simulating algorithm $\SIMEval$ to simulate generating honest players' shares and commitments.
The details of $\SIMEval$ depend on the specifics of the VPSS construction.

The adversary outputs $(b', \vrfin^*, \coalition^*, \{ \vrfpubout_j \}_{j \in \corrupt}),$
where $b'$ is the adversary's guess for the value of $b$,
and $(\vrfin^*, \coalition^*, \{ \vrfpubout_j \}_{j \in \corrupt})$ corresponds to one of the sessions the adversary initiated with $\oracle^{\generate}$.

If $b=1$,
the environment checks if $(\vrfin^*, \coalition^*, \{ \vrfpubout_j \}_{j \in \corrupt})$ corresponds to an  existing session and is consistent with the honest parties' commitments,
outputting a random coin if it does not.
Otherwise, the environment next checks that the aggregated commitment $\dvrf.\combinepub(t, \minparticipants, \coalition^*, \{ \vrfpubout_k \}_{k \in \coalition^*})$ is equal to the simulated commitment $\vrfpubout$ for that session,
if the check does not hold,
the environment outputs $1$.

Otherwise, the environment checks if the adversary's guess $b'$ is equal to $b$;
outputting $1$ if the check holds,
otherwise, outputting $0$.

We define pseudorandomness for a VPSS more formally in Definition~\ref{def:pseudo}.

\begin{definition} \label{def:pseudo}
Let the advantage of an adversary $\adv$ playing the pseudorandomness game as
defined in Figure~\ref{fig:dcvrf-pseudo} be as follows:

\[
  \advant{psdr}{\adv,\dvrf}(\secp, n, t, \minparticipants)
  = \bigabs{\Pr[ \PseudoExp_{\adv,\dvrf}(\secp, n, t, \minparticipants) = 1]  - 1/2}
\]

  A VPSS $\dvrf$ is \defn{pseudorandom} if for all PPT adversaries $\adv$,
  $\advant{psdr}{\adv,\dvrf}$ is a negligible function of $\secp$,
  for $n, t, \minparticipants \in \mathbb{N}$,
  $t \leq \minparticipants \leq n$.
\end{definition}


